\addcontentsline{toc}{subsubsection}{Solution 15. FEOO 2020 Shortlist -- Algebra A1}

\begin{dlsummary}[title={\bfseries Solution 15. FEOO 2020 Shortlist -- Algebra A1},breakable]

\begin{proof}
Let $P(x,y)$ denote the given inequality.

\medskip
\noindent
\textbf{Step 1. Determine $f(1)$.}

From $P(1,1)$,
\[
2f(1)\geq1+f(1)^2,
\]
which gives
\[
(f(1)-1)^2\leq0.
\]
Thus
\[
\boxed{f(1)=1.} \tag{1}
\]

\medskip
\noindent
\textbf{Step 2. Establish a lower bound.}

By $P(x,1)$ and (1),
\[
2xf(x)\geq x^2+xf(x).
\]
Hence
\[
\boxed{xf(x)\geq x^2
\qquad(x\in\mathbb R).} \tag{2}
\]

\medskip
\noindent
\textbf{Step 3. Force equality.}

For $x\ne0$, applying $P(x,1/x)$ gives
\[
2f(1)\geq
1+xf(x)f\left(\frac1{x^2}\right).
\]
By (2) at $1/x^2$, we have
$f(1/x^2)\geq1/x^2>0$. Thus multiplying $xf(x)\geq x^2$ by this
positive value preserves the inequality. Using (1) and (2),
\[
\begin{aligned}
1
&\geq xf(x)f\left(\frac1{x^2}\right)\\
&\geq x^2f\left(\frac1{x^2}\right)\\
&\geq x^2\cdot\frac1{x^2}=1.
\end{aligned}
\]

All inequalities are equalities.
In particular,
\[
xf(x)=x^2.
\]
Therefore,
\[
\boxed{f(x)=x\qquad(x\ne0).} \tag{3}
\]

\medskip
\noindent
\textbf{Step 4. Determine $f(0)$.}

Setting $y=0$ gives
\[
0\geq xf(x)f(0).
\]
For $x\ne0$, by (3),
\[
0\geq x^2f(0).
\]
Consequently,
\[
\boxed{f(0)\leq0.}
\]

\medskip
\noindent
\textbf{Verification.}

Let
\[
f(x)=
\begin{cases}
x,&x\ne0,\\
c,&x=0,
\end{cases}
\qquad c\leq0.
\]

If $xy\ne0$, both sides of the original
inequality equal $2x^2y^2$.

If $x=0$, both sides vanish.

If $y=0$, the inequality becomes
\[
0\geq x^2c,
\]
which holds because $c\leq0$.

Thus all such functions satisfy the
given inequality.
\end{proof}

\end{dlsummary}